\section{Transition energy crossing for polarized protons}
\label{sec:gamma-transition}

\todo[inline]{Draft: this chapter is incomplete and will be extended.}

Collider experiments impose requirements on the luminosity, which, among other
things, limit the longitudinal emittance of the final bunch. During
acceleration, the beam must both cross the transition energy and be split
into 22 bunches by RF gymnastics without blowing up the longitudinal phase
volume.

A possible transition jump scheme for NICA has been considered. Its
characteristic values are the jump amplitude $\Delta\gamma_{tr} = 0.09$ and the
rate of change of the transition energy
$d\gamma_{tr}/dt = \SI{8.5}{s^{-1}}$. With a harmonic RF system, the proposed
jump has little effect on the longitudinal dynamics, since the jump is small
and slow compared with the acceleration rate. With a barrier-bucket RF system,
the jump amplitude is limited by the microwave instability threshold for the
number of particles in a uniform bunch. This limit does not allow reaching
$1\times10^{12}$ particles in the final bunch required for the maximum
luminosity.

The material of this section has been reported
in~\cite{Kolokolchikov:2024hxf,Kolokolchikov:2022akm}.
